\documentclass[10pt]{article}
\usepackage[T1]{fontenc}
\usepackage[utf8]{inputenc}
\usepackage[french]{babel}
\usepackage{amsmath,amssymb,mathrsfs,mathpazo}
\usepackage{geometry}
\geometry{a4paper,margin=2cm}
\usepackage{graphicx,array,tabularx,tabularray,colortbl,multirow,multicol,booktabs,enumitem,pifont,ulem,fancybox,diagbox,cancel}
\usepackage{pstricks,pstricks-add,pst-plot,pst-tree,pst-node,pst-text,pst-func,pst-math,pst-bspline,pst-3dplot,pst-eucl,pst-solides3d}
\let\pscalculate\undefined
\usepackage{pst-bezier}
\usepackage{tikz,pgfplots,tkz-tab}
\usetikzlibrary{arrows,arrows.meta,patterns,calc,positioning,trees,decorations.markings,shapes.geometric,shapes.misc,shapes.arrows,intersections,decorations.pathmorphing,decorations.pathreplacing,backgrounds,shadows,babel}
\pgfplotsset{compat=1.18}
\usepgfplotslibrary{fillbetween}
\usepackage[np]{numprint}
\usepackage{listings,makecell,pas-tableur,twoopt}
\let\barre\undefined
\usepackage[pstricks]{tablvar}
\definecolor{darkgreen}{RGB}{0,100,0}
\providecommand{\red}{\color{red}}
\providecommand{\blue}{\color{blue}}
\providecommand{\green}{\color{green}}
\providecommand{\black}{\color{black}}
\providecommand{\white}{\color{white}}
\providecommand{\gray}{\color{gray}}
\providecommand{\cyan}{\color{cyan}}
\providecommand{\magenta}{\color{magenta}}
\providecommand{\yellow}{\color{yellow}}
\providecommand{\violet}{\color{violet}}
\providecommand{\orange}{\color{orange}}
\providecommand{\euro}{\texteuro}
\providecommand{\EUR}{\texteuro}
\providecommand{\ieme}{\textsuperscript{e}}
\providecommand{\degres}{\ensuremath{^\circ}}
\providecommand{\degree}{\ensuremath{^\circ}}
\providecommand{\vv}[1]{\overrightarrow{#1}}
\providecommand{\Coord}[1]{#1}
\providecommand{\decosix}{}
\providecommand{\starredbullet}{\ensuremath{\bullet}}
\pagestyle{empty}
\setlength{\parindent}{0pt}
\usepackage{adjustbox}\DeclareUnicodeCharacter{2236}{:}\DeclareUnicodeCharacter{2190}{\ensuremath{\leftarrow}}
\providecommand{\R}{\mathbb{R}}
\providecommand{\N}{\mathbb{N}}
\providecommand{\Z}{\mathbb{Z}}
\providecommand{\Q}{\mathbb{Q}}
\providecommand{\D}{\mathbb{D}}
\providecommand{\C}{\mathbb{C}}
\providecommand{\e}{{\mathrm{e}}}
\providecommand{\vect}[1]{\overrightarrow{#1}}
\providecommand{\vectt}[1]{\overrightarrow{\mathrm{#1}}}
\providecommand{\vvt}[1]{\overrightarrow{\mathrm{#1}}}
\providecommand{\np}[2][]{{#2~#1}}
\providecommand{\vv}[1]{{\overrightarrow{#1}}}
\renewcommand{\d}{{\mathrm{d}}}
\providecommand{\degres}{{^\circ}}
\providecommand{\degre}{{^\circ}}
\providecommand{\blue}{\color{blue}}
\providecommand{\red}{\color{red}}
\providecommand{\enskip}{\quad}
\providecommand{\llbracket}{\lbrack\!\lbrack}
\providecommand{\rrbracket}{\rbrack\!\rbrack}
\providecommand{\euro}{{\text{€}}}
\providecommand{\no}{{\text{n°}}}
\providecommand{\interval}[2]{{[#1;#2]}}
\providecommand{\intervFO}[2]{{[#1;#2[}}
\providecommand{\dotfill}{\ldots}

\providecommand{\ds}{\displaystyle}
\providecommand{\textding}[1]{#1}
\providecommand{\Oij}{\(\mathrm{O};\vec{\imath},\vec{\jmath}\)}
\providecommand{\Oijk}{\(\mathrm{O};\vec{\imath},\vec{\jmath},\vec{k}\)}
\providecommand{\Ouv}{\(\mathrm{O};\vec{u},\vec{v}\)}
\providecommand{\intervFF}[2]{\ensuremath{[#1;#2]}}
\renewcommand{\euro}{\ifmmode\text{\texteuro}\else\texteuro\fi}
\DeclareUnicodeCharacter{00B0}{\ifmmode^{\circ}\else\textdegree\fi}
\DeclareUnicodeCharacter{20AC}{\ifmmode\text{\texteuro}\else\texteuro\fi}
% Énoncés : Éducation nationale. Transcription et corrigés : APMEP, Denis Vergès et contributeurs. Découpage et classement : CoopMaths, Sébastien Lozano, Éric Elter et contributeurs. Adaptation de la présentation : Maths Explorer.
% CC BY-SA : https://coopmaths.fr/www/about/
% https://coopmaths.fr/alea/static/bac/2025/tex/bac_2025_11_sujet2_caledonie_1_cor.tex
\begin{document}
\noindent\textit{APMEP / CoopMaths. CC BY-SA. Présentation adaptée par Maths Explorer.}\par\medskip
\smallskip

\begin{enumerate}
	\item Dans un cube toutes les faces sont des carrés de même côté, donc les diagonales de chaque face ont même longueur.

Or dans le triangle HAC, chaque côté est une diagonale d'une face ([HA] est une diagonale de ADHE, [AC] est une diagonale de ABCD et [CH] est une diagonale de CDHG). Les trois côtés du triangle HAC ont donc même longueur $\sqrt2$
, c'est donc un triangle équilatéral.

Or les angles aux sommets d'un triangle équilatéral mesurent 60\ensuremath{^\circ}, donc HAC n'est pas rectangle.

\smallskip

L'affirmation 1 est \textbf{FAUSSE}.
	\item Les droites (DH) et (BF) sont des arêtes opposées du cube, elles sont donc parallèles, et donc coplanaires. Donc (BF) est une droite du plan (DFH).

Comme I est un point de (BF) il appartient au plan (DFH). Donc les droites (DI) et (HF) sont coplanaires.

Puisque (BD) est la droite parallèle à (HF) passant par D (arêtes opposées), et que B et I sont distincts, d'après l'axiome d'Euclide, la droite (DI) n'est pas parallèle à (HF). Or si deux droites sont coplanaires, elles sont soit parallèles soit sécantes. Donc (DI) et (HF) sont sécantes.

\smallskip
	
L'affirmation 2 est \textbf{VRAIE}.
	\item Si le vecteur $\vect{u}$ est normal au plan (FAC) alors il est orthogonal à tout vecteur de ce plan.

On calcule donc le produit scalaire $\vect{u} \cdot \vect{\text{FC}}$, puisque deux vecteurs sont orthogonaux si, et seulement si, leur produit scalaire est nul.
	
	$\vectt{FC} = \begin{pmatrix}1-1\\1-0\\0+1\end{pmatrix} = \begin{pmatrix}0\\1\\-1\end{pmatrix}$
\begin{eqnarray*}
\vect{u} \cdot \vectt{FC} &=& \sin(\alpha) \times 0 + \sin(\pi - \alpha) \times 1 + \sin(-\alpha) \times (-1)\\
		&=& \sin(\pi - \alpha) - \sin(-\alpha)\\
		&=& \sin(\alpha) + \sin(\alpha)\\
		&=& \boxed{2\sin(\alpha) \neq 0} \quad \text{puisque} \:\alpha \ne 0
	\end{eqnarray*}
	Le vecteur $\vect{u}$ n'est donc pas orthogonal au vecteur $\vect{\text{FC}}$. Donc, puisqu'il existe un vecteur de (FAC) qui n'est pas orthogonal à $\vect{u}$, $\vect{u}$ n'est pas un vecteur normal à (FAC).

\smallskip
	
L'affirmation 3 est \textbf{FAUSSE}.
	\item Un cube possède \textbf{8 sommets}.
	
	Pour former un segment, on doit choisir 2 sommets distincts parmi ces 8 sommets. L'ordre n'a pas d'importance (le segment reliant A à B est le même que celui reliant B à A).
	
	On utilise donc une combinaison :
	\[\binom{8}{2} = \dfrac{8\,!}{2\,!(8-2)\,!} = \dfrac{8 \times 7}{2 \times 1} = \dfrac{56}{2} = 28.\]
	Il existe \textbf{28 segments} pouvant relier deux sommets distincts d'un cube.

	\smallskip
	
	L'affirmation 4 est \textbf{FAUSSE}.
\end{enumerate}


\bigskip

\end{document}
